\documentclass{article}
\usepackage[paperwidth=60cm,paperheight=120cm,margin=0pt]{geometry}
\usepackage{fontspec}\setmainfont{Latin Modern Roman}
\usepackage{xeCJK}\setCJKmainfont{Microsoft YaHei}
\usepackage{amsmath, amssymb, siunitx}
\usepackage[american]{circuitikz}
\ctikzset{bipoles/length=1.2cm}
% 不用 standalone：它自家的裁边在 xelatex 下不生效，长图会被 letter 页面截掉尾部。
% 改用 preview 把页面尺寸严格设成内容尺寸（自动裁边）。
\usepackage[active,tightpage]{preview}
\PreviewEnvironment{circuitikz}
\setlength\PreviewBorder{8pt}
\pagestyle{empty}
\begin{document}

\begin{circuitikz}
\small
\node[anchor=east, font=\bfseries\small] at (4.55,0.00) {MCU\ /\ 电源};
\node[anchor=west, font=\bfseries\small] at (4.95,0.00) {网络名};
\node[anchor=west, font=\bfseries\small] at (12.45,0.00) {连接器针脚};
\node[anchor=west, font=\bfseries\small] at (12.30,-0.60) {感为八路灰度（16P，★主方案）};
\node[anchor=west, font=\scriptsize\itshape] at (12.30,-0.98) {I2C 模式：可读数字量+模拟量；5V 供电 33mA；上拉用感为板载 5V/10k 跳线（拔掉 PCB 侧上拉）};
\node[anchor=east, font=\small] at (4.55,-1.84) {GND\ \textperiodcentered\ ①电源负};
\draw (4.85,-1.84) -- (12.30,-1.84);
\node[anchor=south, font=\scriptsize\ttfamily] at (8.57,-1.79) {GND};
\node[anchor=west, font=\small] at (12.48,-1.84) {1\,--\,GND};
\node[anchor=east, font=\small] at (4.55,-2.58) {5V\ \textperiodcentered\ ②电源正 33mA};
\draw (4.85,-2.58) -- (12.30,-2.58);
\node[anchor=south, font=\scriptsize\ttfamily] at (8.57,-2.53) {5V};
\node[anchor=west, font=\small] at (12.48,-2.58) {2\,--\,+5V};
\node[anchor=east, font=\small] at (4.55,-3.32) {KEY\ \textperiodcentered\ 按键（悬空）};
\draw (4.85,-3.32) -- (12.30,-3.32);
\node[anchor=south, font=\scriptsize\ttfamily] at (8.57,-3.27) {GRAY\_KEY};
\node[anchor=west, font=\small] at (12.48,-3.32) {3\,--\,KEY};
\node[anchor=east, font=\small] at (4.55,-4.06) {—\ \textperiodcentered\ 错误输出（可选接 PA11 观察）};
\draw (4.85,-4.06) -- (12.30,-4.06);
\node[anchor=south, font=\scriptsize\ttfamily] at (8.57,-4.01) {GRAY\_ERR};
\node[anchor=west, font=\small] at (12.48,-4.06) {4\,--\,Err};
\node[anchor=east, font=\small] at (4.55,-4.80) {—\ \textperiodcentered\ 并行数据 1~8：本设计不用};
\draw (4.85,-4.80) -- (12.30,-4.80);
\node[anchor=south, font=\scriptsize\ttfamily] at (8.57,-4.75) {NC};
\node[anchor=west, font=\small] at (12.48,-4.80) {5-12\,--\,1 ~ 8};
\node[anchor=east, font=\small] at (4.55,-5.54) {—\ \textperiodcentered\ 串行备份：留排针可跳线};
\draw (4.85,-5.54) -- (12.30,-5.54);
\node[anchor=south, font=\scriptsize\ttfamily] at (8.57,-5.49) {GRAY\_CLK};
\node[anchor=west, font=\small] at (12.48,-5.54) {13\,--\,CLK};
\node[anchor=east, font=\small] at (4.55,-6.28) {—\ \textperiodcentered\ 串行备份：留排针可跳线};
\draw (4.85,-6.28) -- (12.30,-6.28);
\node[anchor=south, font=\scriptsize\ttfamily] at (8.57,-6.23) {GRAY\_DAT};
\node[anchor=west, font=\small] at (12.48,-6.28) {14\,--\,DAT};
\node[anchor=east, font=\small] at (4.55,-7.02) {PB8\ \textperiodcentered\ I2C1\_SCL 重映射};
\draw (4.85,-7.02) -- (12.30,-7.02);
\node[anchor=south, font=\scriptsize\ttfamily] at (8.57,-6.97) {I2C\_GRAY\_SCL};
\node[anchor=west, font=\small] at (12.48,-7.02) {15\,--\,SCL};
\node[anchor=east, font=\small] at (4.55,-7.76) {PB9\ \textperiodcentered\ I2C1\_SDA 重映射};
\draw (4.85,-7.76) -- (12.30,-7.76);
\node[anchor=south, font=\scriptsize\ttfamily] at (8.57,-7.71) {I2C\_GRAY\_SDA};
\node[anchor=west, font=\small] at (12.48,-7.76) {16\,--\,SDA};
\draw[rounded corners=2pt] (12.30,-1.42) rectangle (18.40,-8.06);
\node[anchor=west, font=\bfseries\small] at (12.30,-8.66) {亚博八路灰度（16P，备用）};
\node[anchor=west, font=\scriptsize\itshape] at (12.30,-9.04) {针 1/2 与感为相反（5V 在 1 脚）→ 必须独立座子，不可和感为共用；同样支持 I2C};
\node[anchor=east, font=\small] at (4.55,-9.90) {5V\ \textperiodcentered\ ①电源};
\draw (4.85,-9.90) -- (12.30,-9.90);
\node[anchor=south, font=\scriptsize\ttfamily] at (8.57,-9.85) {5V};
\node[anchor=west, font=\small] at (12.48,-9.90) {1\,--\,5V};
\node[anchor=east, font=\small] at (4.55,-10.64) {GND\ \textperiodcentered\ ②接地};
\draw (4.85,-10.64) -- (12.30,-10.64);
\node[anchor=south, font=\scriptsize\ttfamily] at (8.57,-10.59) {GND};
\node[anchor=west, font=\small] at (12.48,-10.64) {2\,--\,GND};
\node[anchor=east, font=\small] at (4.55,-11.38) {KEY\ \textperiodcentered\ 按键（悬空）};
\draw (4.85,-11.38) -- (12.30,-11.38);
\node[anchor=south, font=\scriptsize\ttfamily] at (8.57,-11.33) {GRAY\_KEY};
\node[anchor=west, font=\small] at (12.48,-11.38) {3\,--\,KEY};
\node[anchor=east, font=\small] at (4.55,-12.12) {—\ \textperiodcentered\ 复位（悬空）};
\draw (4.85,-12.12) -- (12.30,-12.12);
\node[anchor=south, font=\scriptsize\ttfamily] at (8.57,-12.07) {NC};
\node[anchor=west, font=\small] at (12.48,-12.12) {4\,--\,RESET};
\node[anchor=east, font=\small] at (4.55,-12.86) {—\ \textperiodcentered\ 并行 IO X1~X8：本设计不用};
\draw (4.85,-12.86) -- (12.30,-12.86);
\node[anchor=south, font=\scriptsize\ttfamily] at (8.57,-12.81) {NC};
\node[anchor=west, font=\small] at (12.48,-12.86) {5-12\,--\,X1 ~ X8};
\node[anchor=east, font=\small] at (4.55,-13.60) {—\ \textperiodcentered\ 串口 TX/RX：留排针可跳线};
\draw (4.85,-13.60) -- (12.30,-13.60);
\node[anchor=south, font=\scriptsize\ttfamily] at (8.57,-13.55) {GRAY\_TX};
\node[anchor=west, font=\small] at (12.48,-13.60) {13\,--\,TX};
\node[anchor=east, font=\small] at (4.55,-14.34) {—\ \textperiodcentered\ 串口 TX/RX：留排针可跳线};
\draw (4.85,-14.34) -- (12.30,-14.34);
\node[anchor=south, font=\scriptsize\ttfamily] at (8.57,-14.29) {GRAY\_RX};
\node[anchor=west, font=\small] at (12.48,-14.34) {14\,--\,RX};
\node[anchor=east, font=\small] at (4.55,-15.08) {PB8\ \textperiodcentered\ I2C1\_SCL 重映射};
\draw (4.85,-15.08) -- (12.30,-15.08);
\node[anchor=south, font=\scriptsize\ttfamily] at (8.57,-15.03) {I2C\_GRAY\_SCL};
\node[anchor=west, font=\small] at (12.48,-15.08) {15\,--\,SCL};
\node[anchor=east, font=\small] at (4.55,-15.82) {PB9\ \textperiodcentered\ I2C1\_SDA 重映射};
\draw (4.85,-15.82) -- (12.30,-15.82);
\node[anchor=south, font=\scriptsize\ttfamily] at (8.57,-15.77) {I2C\_GRAY\_SDA};
\node[anchor=west, font=\small] at (12.48,-15.82) {16\,--\,SDA};
\draw[rounded corners=2pt] (12.30,-9.48) rectangle (18.40,-16.12);
\node[anchor=west, font=\bfseries\small] at (12.30,-16.72) {超声波 1（左，XH2.54-4P）};
\node[anchor=west, font=\scriptsize\itshape] at (12.30,-17.10) {两模块 Trig 独立、软件分时触发（防串扰）；Echo 5V→3.1V 分压};
\node[anchor=east, font=\small] at (4.55,-17.96) {5V};
\draw (4.85,-17.96) -- (12.30,-17.96);
\node[anchor=south, font=\scriptsize\ttfamily] at (8.57,-17.91) {5V};
\node[anchor=west, font=\small] at (12.48,-17.96) {1\,--\,VCC};
\node[anchor=east, font=\small] at (4.55,-18.70) {PA4};
\draw (4.85,-18.70) -- (12.30,-18.70);
\node[anchor=south, font=\scriptsize\ttfamily] at (8.57,-18.65) {US\_L\_TRIG};
\node[anchor=west, font=\small] at (12.48,-18.70) {2\,--\,Trig};
\node[anchor=east, font=\small] at (4.55,-19.44) {PA5\ \textperiodcentered\ 2k/3.3k 分压};
\draw (4.85,-19.44) -- (12.30,-19.44);
\node[anchor=south, font=\scriptsize\ttfamily] at (8.57,-19.39) {US\_L\_ECHO};
\node[anchor=west, font=\small] at (12.48,-19.44) {3\,--\,Echo};
\node[anchor=east, font=\small] at (4.55,-20.18) {GND};
\draw (4.85,-20.18) -- (12.30,-20.18);
\node[anchor=south, font=\scriptsize\ttfamily] at (8.57,-20.13) {GND};
\node[anchor=west, font=\small] at (12.48,-20.18) {4\,--\,GND};
\draw[rounded corners=2pt] (12.30,-17.54) rectangle (18.40,-20.48);
\node[anchor=west, font=\bfseries\small] at (12.30,-21.08) {超声波 2（右，XH2.54-4P）};
\node[anchor=east, font=\small] at (4.55,-21.92) {5V};
\draw (4.85,-21.92) -- (12.30,-21.92);
\node[anchor=south, font=\scriptsize\ttfamily] at (8.57,-21.87) {5V};
\node[anchor=west, font=\small] at (12.48,-21.92) {1\,--\,VCC};
\node[anchor=east, font=\small] at (4.55,-22.66) {PA8};
\draw (4.85,-22.66) -- (12.30,-22.66);
\node[anchor=south, font=\scriptsize\ttfamily] at (8.57,-22.61) {US\_R\_TRIG};
\node[anchor=west, font=\small] at (12.48,-22.66) {2\,--\,Trig};
\node[anchor=east, font=\small] at (4.55,-23.40) {PA15\ \textperiodcentered\ 关 JTAG};
\draw (4.85,-23.40) -- (12.30,-23.40);
\node[anchor=south, font=\scriptsize\ttfamily] at (8.57,-23.35) {US\_R\_ECHO};
\node[anchor=west, font=\small] at (12.48,-23.40) {3\,--\,Echo};
\node[anchor=east, font=\small] at (4.55,-24.14) {GND};
\draw (4.85,-24.14) -- (12.30,-24.14);
\node[anchor=south, font=\scriptsize\ttfamily] at (8.57,-24.09) {GND};
\node[anchor=west, font=\small] at (12.48,-24.14) {4\,--\,GND};
\draw[rounded corners=2pt] (12.30,-21.50) rectangle (18.40,-24.44);
\node[anchor=west, font=\bfseries\small] at (12.30,-25.04) {OLED 显示屏（XH2.54-4P）};
\node[anchor=west, font=\scriptsize\itshape] at (12.30,-25.42) {走 I2C2（3.3V 上拉）；与 MPU6050 同总线、与灰度总线完全隔离};
\node[anchor=east, font=\small] at (4.55,-26.28) {3V3};
\draw (4.85,-26.28) -- (12.30,-26.28);
\node[anchor=south, font=\scriptsize\ttfamily] at (8.57,-26.23) {3V3};
\node[anchor=west, font=\small] at (12.48,-26.28) {1\,--\,VCC};
\node[anchor=east, font=\small] at (4.55,-27.02) {GND};
\draw (4.85,-27.02) -- (12.30,-27.02);
\node[anchor=south, font=\scriptsize\ttfamily] at (8.57,-26.97) {GND};
\node[anchor=west, font=\small] at (12.48,-27.02) {2\,--\,GND};
\node[anchor=east, font=\small] at (4.55,-27.76) {PB10\ \textperiodcentered\ I2C2\_SCL};
\draw (4.85,-27.76) -- (12.30,-27.76);
\node[anchor=south, font=\scriptsize\ttfamily] at (8.57,-27.71) {I2C\_AUX\_SCL};
\node[anchor=west, font=\small] at (12.48,-27.76) {3\,--\,SCL};
\node[anchor=east, font=\small] at (4.55,-28.50) {PB11\ \textperiodcentered\ I2C2\_SDA};
\draw (4.85,-28.50) -- (12.30,-28.50);
\node[anchor=south, font=\scriptsize\ttfamily] at (8.57,-28.45) {I2C\_AUX\_SDA};
\node[anchor=west, font=\small] at (12.48,-28.50) {4\,--\,SDA};
\draw[rounded corners=2pt] (12.30,-25.86) rectangle (18.40,-28.80);
\node[anchor=west, font=\bfseries\small] at (12.30,-29.40) {MPU6050 六轴（XH2.54-4P）};
\node[anchor=west, font=\scriptsize\itshape] at (12.30,-29.78) {备用，暂不参与控制；AD0 接地 → 地址 0x68};
\node[anchor=east, font=\small] at (4.55,-30.64) {3V3};
\draw (4.85,-30.64) -- (12.30,-30.64);
\node[anchor=south, font=\scriptsize\ttfamily] at (8.57,-30.59) {3V3};
\node[anchor=west, font=\small] at (12.48,-30.64) {1\,--\,VCC};
\node[anchor=east, font=\small] at (4.55,-31.38) {GND};
\draw (4.85,-31.38) -- (12.30,-31.38);
\node[anchor=south, font=\scriptsize\ttfamily] at (8.57,-31.33) {GND};
\node[anchor=west, font=\small] at (12.48,-31.38) {2\,--\,GND};
\node[anchor=east, font=\small] at (4.55,-32.12) {PB10\ \textperiodcentered\ I2C2\_SCL};
\draw (4.85,-32.12) -- (12.30,-32.12);
\node[anchor=south, font=\scriptsize\ttfamily] at (8.57,-32.07) {I2C\_AUX\_SCL};
\node[anchor=west, font=\small] at (12.48,-32.12) {3\,--\,SCL};
\node[anchor=east, font=\small] at (4.55,-32.86) {PB11\ \textperiodcentered\ I2C2\_SDA};
\draw (4.85,-32.86) -- (12.30,-32.86);
\node[anchor=south, font=\scriptsize\ttfamily] at (8.57,-32.81) {I2C\_AUX\_SDA};
\node[anchor=west, font=\small] at (12.48,-32.86) {4\,--\,SDA};
\draw[rounded corners=2pt] (12.30,-30.22) rectangle (18.40,-33.16);
\node[anchor=west, font=\bfseries\small] at (12.30,-33.76) {12V 电压检测（分压）};
\node[anchor=west, font=\scriptsize\itshape] at (12.30,-34.14) {100k / 10k 分压 + 100nF；只做低压告警，不参与控制};
\node[anchor=east, font=\small] at (4.55,-35.00) {12V\ \textperiodcentered\ 经 100k};
\draw (4.85,-35.00) -- (12.30,-35.00);
\node[anchor=south, font=\scriptsize\ttfamily] at (8.57,-34.95) {VBAT\_SENSE};
\node[anchor=west, font=\small] at (12.48,-35.00) {—\,--\,分压中点 → PB0 (ADC1\_IN8)};
\draw[rounded corners=2pt] (12.30,-34.58) rectangle (18.40,-35.30);
\node[anchor=west, font=\scriptsize\itshape] at (0,-36.20) {网络名 = 嘉立创 EDA 里的网络标号：两边同名即同一根线。电源（5V / 3V3 / GND / 12V）同样用网络标号表示，不画共用母线。逐针定义、电气约束与改动规则见 电气连接规范.md；引脚契约见 docs/引脚分配表.md。};
\end{circuitikz}
\end{document}
